% from main.tex
%%%%%%%%%%  scan idx 566  |  folio 553  %%%%%%%%%%
modules localement libres de rang \(n_1,\ldots,n_r\)) par recollement du \(\mathcal O\)-module libre \(T(\mathcal O^{n_1},\ldots,\mathcal O^{n_r})\) sur des cartes locales trivialisantes moyennant la représentation \(\alpha\). De façon précise, à la suite \(E_1,\ldots,E_r\) correspond un \(G\)-torseur à droite \(P=P_{E_1,\ldots,E_r}\) tel que l'on a un isomorphisme canonique
\begin{equation}\tag{6.1.3}
 E_i\xrightarrow{\sim}P\overset{G}{\wedge}\mathcal O^{n_i}\qquad(i=1,\ldots,r),
\end{equation}
où \(\overset{G}{\wedge}\) désigne le produit contracté, et \(G\) opère sur \(\mathcal O^{n_i}\) via la \(i\)-ième projection \(G\to\Gl(n_i)\) ; alors toute représentation
\[
 \alpha:G\longrightarrow\Aut_{\mathcal O}(M)
\]
fait opérer \(G\) sur \(\mathcal O^m\) par des \(\mathcal O\)-automorphismes, et définit ainsi un multifoncteur \(T=T_\alpha\) tel que
\begin{equation}\tag{6.1.4}
 T_\alpha(E_1,\ldots,E_r)=P_{E_1,\ldots,E_r}\overset{G}{\wedge} M.
\end{equation}
De plus, pour toute représentation composée \(\alpha=\alpha_1\circ\alpha_2\) on trouve
\begin{equation}\tag{6.1.5}
 T_\alpha=T_{\alpha_1}\circ T_{\alpha_2}.
\end{equation}

\subsection*{6.2}
Pour tout \(\mathcal O\)-module localement libre \(E\) de rang \(n\) nous désignerons par \(\det(E)\) le \(\mathcal O\)-module inversible \(\bigwedge\limits^n E\). Gardant les notations de 6.1, le foncteur \(\det(T(E_1,\ldots,E_r))\) est déterminé, à isomorphisme canonique près, par la représentation correspondante
\begin{equation}\tag{6.2.1}
 \det\circ\alpha_T:G\longrightarrow\Gm=\Gl(1).
\end{equation}

Supposons maintenant que \(X\) soit le topos annelé des faisceaux sur le site annelé \((\Sch/S)\) muni de la topologie étale. Alors on peut calculer \(\det(T(E_1,\ldots,E_r))\), ou ce qui revient au même, la représentation \(\det\circ\alpha_T\), à l'aide de la

